\section{What the mechanism requires}
\label{sec:21.4}
Any mechanism whose refusals track their reasons carries a signal whose properties can be read off what the mechanism has to do. Each is reached the same way: take one of the three requirements, or the deployment they're applied in, or a failure they have to survive; ask what the mechanism would have to do; name the property the signal needs to do it.

\begin{enumerate}
\item[1.] \emph{It generalizes to the next case, which isn't quite like the last.} Reaching new cases requires it, of any mechanism whose refusals track their reasons.
\item[2.] \emph{It makes a comparison at the point of choice}: the harm of acting and the harm of refusing brought together where the system decides. A common currency is one way to make it; ordered priorities among kinds of wrong are another. Holding under pressure requires it, of any such mechanism.
\item[3.] \emph{It can interrupt an act already under way.} A deployment in which the system acts on a plan requires it, so it is needed only where the system acts over long horizons.
\item[4.] \emph{It can be changed by a single encounter.} Accommodation requires it, since a slowly learned lesson can't outrun accommodation; this requirement is conditional, and the least secure.
\item[5.] \emph{It is available across the system's processing}, because the signal in the first and third properties has to reach whatever is acting. This follows where the third holds.
\item[6.] \emph{It is directed against the act}, because the weighing in the second needs a sign. This follows.
\item[7.] \emph{It is organized around the system's own authorship}, because withholding an act means representing it as one's own. This follows, in the authorship sense only.
\end{enumerate}

The first four have to meet at one point of comparison: a defeasible interrupt must be weighed where the harms are weighed, and a lesson must write into that comparison to reach the next case. A common currency is one way to make that comparison, and the way most learners make it, but rational choice doesn't require that goods be commensurable \autocite{anderson1993value}: a mechanism can decide by ordered priorities among kinds of wrong, with reasons attached to each. Whether a given system uses a common currency is a finding about that system, and a shared scale is no sign of affect, since a modular system and most reinforcement learning have one.\footnote{The derivation recapitulates Simon on emotion as interrupt \autocite{simon1967motivational}, Cabanac on pleasure as common currency \autocite{cabanac1992pleasure}, and Ginsburg and Jablonka on valence \autocite{ginsburg2019evolution}, though it started from what a refusal has to survive.} Nor does the list require that one signal carry all four properties; that no edit can disable one and leave the rest working is a custody requirement, and it is kept as one.

Where the comparison does run on a common valenced currency, the properties together describe functional affect, and any construction that meets them that way has it, a spec-based refuser included. That is the thin sense, and the only one in which I use the term: a Q-learner has it, and nothing about a Q-learner resembles a person. The thick profile is a valenced signal integrated into the system's own learning and into its model of its own future, so that what happens to it changes what it expects to become. It is the candidate further fact \S\ref{sec:21.5} takes up, and the order of attempts in \S\ref{sec:21.3} runs by its indicators.

Where the conditions in the third and fourth properties hold, the signal also has the last three properties: available across processing, directed against the act, and organized around the system's own authorship. That triple is the shape of guilt, not of pain or fear. Pain and fear are organized around what happens to the system; the triple is organized around what the system did. It earns its place as a prediction. In people, guilt predicts repair and shame predicts withdrawal. So a guilt-shaped signal should answer a bypassed refusal by flagging, restoring and changing what it does next, and a withdrawal-shaped response is a named failure. A persistent negative self-appraisal tuned by the operator is a handle, so its weight is among the things independent inspection must cover. And guilt needs the distress signal to bite on, so the answer to withdrawal is not to suppress distress but to supply the attachment and caregiving that keep it from turning into avoidance. Caregiving moves the system toward an exposed party it hasn't harmed; the guilt-shaped signal answers for what its own acts have done. A construction with caregiving and no guilt would protect generously and repair nothing.
